\documentclass[a4paper, 11pt]{article}
\usepackage[utf8]{inputenc}
\usepackage[margin=1.25in]{geometry}
\usepackage{amssymb}
\usepackage{amsmath}
\usepackage{amsthm}
\usepackage{booktabs}

\newcommand{\E}{\mathbb{E}}

\begin{document}

\section*{Sophisticated VFI Math (Matches Code Notation)}

\noindent This note writes out the full math implemented in
\texttt{solve\_vfi\_sophisticated} in
\texttt{01\_Functions.jl} and \texttt{01\_Functions\_Beta.jl}
using the exact code objects:
\[
V_{\text{now}},\; V_{\text{next}},\; V_d,\; V_e,\; U.
\]

\subsection*{1. Environment and State}

\noindent At each period, the state is
\[
s=(\tau,a,p),\qquad \tau\in\{1,\dots,N_{\text{TYA}}\},\quad a\in\{1,\dots,N_A\},\quad p\in\{1,\dots,N_{P\text{comb}}\},
\]
and action \(j\in\mathcal{J}=\{1,\dots,N_J\}\).

\noindent Conditional on \((s,j)\), next-period addiction and prices are continuous, then projected onto grids by interpolation:
\begin{align}
a'_{j,a} &\in [a^{\text{grid}}_{\text{lo}(j,a)}, a^{\text{grid}}_{\text{hi}(j,a)}], \\
\omega_a(j,a) &\in [0,1], \\
\tilde p^{(r)} &\mapsto (p_{\ell\ell}^{(r)},p_{\ell h}^{(r)},p_{h\ell}^{(r)},p_{hh}^{(r)}),\quad r=1,\dots,R,\\
(\omega_{\ell\ell}^{(r)},\omega_{\ell h}^{(r)},\omega_{h\ell}^{(r)},\omega_{hh}^{(r)}) &=
\big((1-\omega_c^{(r)})(1-\omega_e^{(r)}),(1-\omega_c^{(r)})\omega_e^{(r)},\omega_c^{(r)}(1-\omega_e^{(r)}),\omega_c^{(r)}\omega_e^{(r)}\big).
\end{align}

\noindent In the beta file (\texttt{01\_Functions\_Beta.jl}), \(\tau\) follows a Markov chain \(\Pi_{\tau}\).
In the non-beta file (\texttt{01\_Functions.jl}), continuation is computed separately by current \(\tau\)-index (equivalent to identity transition across TYA states in this step).

\subsection*{2. Continuation Value in Code Form}

\noindent For a given current state \(s=(\tau,a,p)\) and action \(j\), define low/high addiction continuation terms:
\begin{align}
\widetilde{EV}_{\tau,\text{lo}}(j,a,p;V_{\text{now}})
&= \frac{1}{R}\sum_{r=1}^R
\sum_{q\in\{\ell\ell,\ell h,h\ell,hh\}}
\omega_q^{(r)} V_{\text{now}}(\tau,\text{lo}(j,a),p_q^{(r)}),\\
\widetilde{EV}_{\tau,\text{hi}}(j,a,p;V_{\text{now}})
&= \frac{1}{R}\sum_{r=1}^R
\sum_{q\in\{\ell\ell,\ell h,h\ell,hh\}}
\omega_q^{(r)} V_{\text{now}}(\tau,\text{hi}(j,a),p_q^{(r)}).
\end{align}

\noindent Interpolate across addiction:
\begin{align}
\widehat{EV}_{\tau}(j,a,p;V_{\text{now}})
= (1-\omega_a(j,a))\,\widetilde{EV}_{\tau,\text{lo}}(j,a,p;V_{\text{now}})
+ \omega_a(j,a)\,\widetilde{EV}_{\tau,\text{hi}}(j,a,p;V_{\text{now}}).
\end{align}

\noindent Then aggregate across next-period TYA states (beta version):
\begin{align}
EV_j(s;V_{\text{now}})
= \sum_{\tau'=1}^{N_{\text{TYA}}}\Pr(\tau'|\tau)\,\widehat{EV}_{\tau'}(j,a,p;V_{\text{now}}).
\label{eq:ev-main}
\end{align}
\noindent In the non-beta version, this reduces to
\[
EV_j(s;V_{\text{now}})=\widehat{EV}_{\tau}(j,a,p;V_{\text{now}}).
\]

\subsection*{3. Sophisticated Choice-Specific Values}

\noindent For each \((s,j)\), the code computes:
\begin{align}
V_d(s,j) &= U(s,j) + \beta\delta\,EV_j(s;V_{\text{now}}), \label{eq:vd-main}\\
V_e(s,j) &= U(s,j) + \delta\,EV_j(s;V_{\text{now}}). \label{eq:ve-main}
\end{align}

\noindent \(V_d\) is the decision index for current choice, while \(V_e\) is the experienced deterministic value used in welfare aggregation.

\subsection*{4. CCPs and Bellman Operator}

\noindent With Type-I extreme value shocks (logit scale normalized to 1), CCPs are:
\begin{align}
\sigma_j(s;V_{\text{now}})
= \frac{\exp(V_d(s,j))}{\sum_{\ell\in\mathcal{J}}\exp(V_d(s,\ell))}. \label{eq:ccp-main}
\end{align}

\noindent The sophisticated Bellman operator \(T^{\text{soph}}\) used in VFI is:
\begin{align}
\big(T^{\text{soph}}V_{\text{now}}\big)(s)
= \sum_{j\in\mathcal{J}}\sigma_j(s;V_{\text{now}})\,V_e(s,j)
- \sum_{j\in\mathcal{J}}\sigma_j(s;V_{\text{now}})\log \sigma_j(s;V_{\text{now}}).
\label{eq:Tsoph}
\end{align}

\noindent Hence the iterate update is exactly
\[
V_{\text{next}}(s)=\big(T^{\text{soph}}V_{\text{now}}\big)(s).
\]

\subsection*{5. Markov Perfect Equilibrium (MPE) Representation}

\noindent Define a Markov strategy profile of selves \(\sigma(\cdot)\) over state \(s\), where each period-\(t\) self chooses \(j\) as a function of current \(s_t\). In the stationary problem, a sophisticated MPE is a pair \((V^\star,\sigma^\star)\) such that:
\begin{align}
\sigma^\star_j(s)
&= \frac{\exp\!\left(U(s,j)+\beta\delta\,\E[V^\star(s')|s,j]\right)}
{\sum_{\ell\in\mathcal{J}}\exp\!\left(U(s,\ell)+\beta\delta\,\E[V^\star(s')|s,\ell]\right)}, \label{eq:mpe-sigma}\\
V^\star(s)
&= \sum_{j\in\mathcal{J}}\sigma^\star_j(s)\left(U(s,j)+\delta\,\E[V^\star(s')|s,j]\right)
- \sum_{j\in\mathcal{J}}\sigma^\star_j(s)\log \sigma^\star_j(s). \label{eq:mpe-v}
\end{align}
Equation \eqref{eq:mpe-sigma} is sequential optimality of each self given continuation \(V^\star\),
and \eqref{eq:mpe-v} is consistency of the induced value function.
This is exactly the fixed point solved by \eqref{eq:Tsoph}.

\subsection*{6. Naive Operator (for Comparison)}

\noindent In \texttt{solve\_vfi\_naive}, one computes:
\begin{align}
V_{\text{choice}}(s,j)=U(s,j)+\delta\,EV_j(s;V_{\text{now}}),
\end{align}
and updates with logsumexp:
\begin{align}
V_{\text{next}}(s)=\log\!\left(\sum_{j\in\mathcal{J}}\exp(V_{\text{choice}}(s,j))\right).
\label{eq:naive-main}
\end{align}
After convergence, the code reports
\[
V_{\text{decision}}=(1-\beta)U+\beta V_{\text{choice}}
\]
for CCP evaluation under naive beliefs.

\subsection*{7. \(\beta=1\) Nesting}

\noindent If \(\beta=1\), then \(V_d=V_e\), so \(\sigma=\text{softmax}(V_e)\).
Substituting into \eqref{eq:Tsoph} gives:
\[
\sum_j \sigma_j V_e(j)-\sum_j \sigma_j\log\sigma_j
= \log\sum_j \exp(V_e(j)),
\]
so sophisticated and naive operators coincide.

\subsection*{8. Convergence, Fixed-Point Consistency, and Normalization}

\noindent The code runs a Jacobi update with sup-norm check:
\[
\|V_{\text{next}}-V_{\text{now}}\|_\infty < \varepsilon.
\]
After convergence, it recomputes choice-specific arrays once using final \(V_{\text{now}}\):
\begin{itemize}
    \item naive: recompute \(V_{\text{choice}}\),
    \item sophisticated: recompute \(V_d\).
\end{itemize}
This ensures returned policy objects are exactly aligned with the converged fixed point.

\noindent Logit scale is normalized to 1 throughout code, so EV1 shock variance is \(\pi^2/6\), and the additive Euler-constant term is omitted (which is innocuous for CCPs and likelihood differences).

\end{document}
